\uuid{wdfE}
\exo7id{4201}
\titre{exo7 4201}
\auteur{quercia}
\organisation{exo7}
\datecreate{2010-03-11}
\isIndication{false}
\isCorrection{true}
\chapitre{Equation différentielle}
\sousChapitre{Equations aux dérivées partielles}
\cpge{mpsi-mp2i.fonctions-deux-variables.b, pcsi.fonctions-deux-variables.b, ptsi.fonctions-deux-variables.b, mpsi-mp2i.primitives-equations-differentielles.b, pcsi.techniques-calcul-primitives.b, ptsi.techniques-calcul-primitives.b, bcpst1.fonctions-deux-variables.c, bcpst1.equations-differentielles-lineaires.a, mp-mpi.calcul-differentiel-optimisation.a}
\module{Analyse}
\niveau{L2}
\difficulte{}

\contenu{
\question{
Résoudre l'équation $2\frac{\partial f}{\partial x} + 3\frac{\partial f}{\partial y} = 4f$ avec la {\it condition aux limites\/}
: $f(t,t) = t\ (t \in \R)$.

(\'Etudier $\varphi : t  \mapsto f(a+bt,a+ct)$ avec $a,b,c$ bien choisis)
}
\reponse{
$f(x,y) = (3x-2y)e^{4(x-y)}$.
}
}
